\section{Lezione 14} \hfill 31/03/2026\\
La douglasia,assieme al pino strobo, è l'unica specie ad aver avuto successo in Italia ed Europa. Oggi ci sono discussioni sull'aumento della clotivazione\\
Originaria dell... montagne rocciose.\\
Albero di primissima grandezza, può raggiungere i 100 metri di altezza, facilmente i 70 metri. Nel suo areale piove tantissimo (Seattle), più di 3000 mm all'anno.\\
Ha un legno che assomiglia a quello del larice, al contrario del larice, quello di douglasia ha un passaggio tra gli anelli più sfumati.\\
È stato famoso in Italia per il suo rigatino. È un legno molto durabile. Molto utilizzata in passato per fare infissi (può permettersi di prendere acqua senza marcire immediatamente).\\
L'ambiente ideale è la faggeta: sono stati fatti tentativi di piantarla più in alto e più in basso.\\
È stata usata molto in Francia quando è arrivato il cancro corticale del castagno: è stato piantato in modo da trovare l'alternativa al castagno. Il 7\% della superficie forestale tedesca è occupata da Douglas; si sta ipotizzando di aumentare il numero al 10.\\
In Italia la douglasia non vive bene : pur essendoci stazioni oceaniche, ci sono spesso problemi di siccità estiva, che rende l'accrescimento e lo sviluppo più difficoltoso. La particella più famosa in Italia è al Vallombrosa.\\
La douglasia è riconoscibile dagli strobili (simili a quello dell'abete rosso, ma con linguetta bifida che escono dal seme) e dagli aghi che assomigliano a quelli del larice (che sanno fortemente di limone).\\
È una specie definitiva.\\
\subsection{Abete rosso - Picea abies - Spruce}
È la specie forestale più importante d'Europa; in Italia appartiene alla triade con il faggio e castagno.\\
Il suo areale è centro-europeo: dai Pirenei fino ai Carpazi e (forse) dall'Appennino Tosco-Emiliano fino alla Svezia Centrale.\\
È una specie, dal punto di vista ecologico, continentale.\\
Specie leader della regione mesalpica, definitiva. Presente anche nella regione esalpica ed endalpica ? \\
Mesofila , gran opportunista. Rispetto alle temperature, ha un gradiente dal piano montano fino al subalpino. Non sopporta molto la siccità estiva (nel 2003 ci fu l'eliminazione degli imboschimenti planiziali e collinari in Italia).\\
Apparato radicale superficiale e fascicolato, che lo rende potenzialmente instabile soprattutto in impianti artificiali (molto presenti).\\
La coltura dell'abete rosso nasce in Germania nell '800, che li piantavano nei campi. La prima crisi è avvenuta negli anni '80, con l'inquinamento atmosferico ed il deperimento dei popolamenti.\\
Esiste anche il Picea morica, attorno ai Carpazi, spesso usato in italia nei giardini. Picea obovata, nella siberia meridionale. Si distinguono non molto, la obovata ha coni e aghi più piccoli.\\
In america c'è il black spruce ed altre picee.\\
Il legno dell'abete rosso è utilizzato per fare tutto: legna da ardere, pellets (nati come uso degli scarti di segheria), tavolame, travi, perline, pavimenti, pannelli di ogni tipo (lamellari), modellismo, e strumenti musicali (abete rosso di risonanza, da fuori è irriconoscibile, gli anelli del tronco hanno delle piegature). Una partita di abete rosso di risonanza è di un albero, con un mercato venduto a chili.\\
È in grado di colonizzare dal piano montano fino ad il subalpino. Il problema è che nella rinnovazione ha comportamenti diversi in funzione della quota in cui si sviluppa:
\begin{itemize}
    \item nel piano montano la rinnovazione è di norma pronta ed abbondante, con una pasciona ogni 4 anni. La rinnovazione dev'essere liberata molto in fretta (avere luce), sennò deperisce e muore;
    \item nel piano sub-alpino la rinnovazione diventa molto lenta, e lo sviluppo è lentissimo (per crescere deve uscire dalla neve).
\end{itemize}
È importante capire dove ci si trova per gestire una pecceta, capendo se il piano è montano, subalpino ed il framezzo. Gli indicatori che ci possono dire dove ci si trova:
\begin{itemize}
    \item abete rosso stesso: nel piano montano ha una tipica forma a pettine: i rami principali sono cadenti ed i rami di secondo ordine penzolano giù. chiome molto larghe
    \item nel piano sub-alpino: la forma dell'albero è molto colonnare ed i rametti sono portati intorno ai rami di primo ordine (a scovolo);
    \item la forma a spazzola , con rami larghi ed i secondari sono sono attorno ai primari, nel montano superiore
    \item struttura del bosco: nel montano i boschi sono chiusi, monoplani (le microstazioni sfavorevoli non sono così troppo sfavorevoli, e la struttura è definita dalla competizione);
    \item nel piano subalpino la struttura è aperta (a cluster, collettivi), le condizioni microstazionali sfavorevoli diventano così tanto sfavorevoli da inibire lo sviluppo. I cluster sono separati, poichè nel mezzo l'abete rosso non riesce a crescere: il popolamento è un mosaico di collettivi. In questi collettivi, di vario numero di piante, è molto importante la facilitazione tra gli individui (superiore alla competizione). Più si sale di quota e più il popolamento si apre.
    \item la vegetazione al suolo: la vegetazione nel piano montano non è mai troppo densa (le piante facendo un popolamento monoplano non porta troppa luce al suolo), tip muschi o erbe a foglia larga. Nel piano subalpino c'è tanta luce e poca densità, molte graminee, poace. 
\end{itemize}
L'abete rosso necessita di almeno 3 mesi di sviluppo vegetativo: non raggiunge il limite del bosco. Un'altra cosa che necessita, soprattutto nel piano subalpino, è di almeno 3 ore di luce diretta (calda) al giorno.\\
Quando trattiamo una pecceta, vogliamo sia raccogliere che indurre la rinnovazione. Ci sono selvicolture diverse a seconda del piano montano o sublalpino.\\
Nei collettivi, piccoli boschetti di poche piante, le chiome arrivano fino a terra.\\
In cima ad un monte la stazione è più favorevole rispetto ad una stazione di fondovalle (diversa neve e luce). Il collettivo rende la stazione negativa ancora più negativa e più positiva la migliore. (....apporto di s.o.).\\
Nel piano subalpino la struttura è a collettivi: è più importante la facilitazione rispetto alla competizione. Occorre dare almeno 3 ore di luce al giorno. L'unico trattamento è il taglio raso del collettivo. Se non viene trattato in questo modo, viene destabilizzato tutto il popolamento. L'eliminazione dei collettivi va studiata, perchè il popolamento compie la funzione di protezione.\\
Tra il piano montano superiore ed il subalpino basso (1600-2000 m), possiamo svolgere delle utilizzazioni con lo scopo di utilizzare il legname. Occorre fare cantieri grandi, facendo tagli a fessura.\\
Negli anni 80-90, il taglio raso è stato quasi bandito. In questo modo è stato provato di fare tagli successivi nel montano superiore: c'è stata la destabilizzazione dei collettivi e l'invasione delle megaformie (suolo freddo e luce diffusa). Per questi motivi sono stati creati dei tagli a fessura: tagli larghi mezza volta- una volta l'altezza delle piante circostanti, e lunghi in base ai capisaldi forti del bosco (margini naturali del bosco molto forti, es. torbiere). I tagli a fessura:
\begin{itemize}
    \item perpendicolari alla pendenza: limitano il movimento della neve (es. valanghe, reptazione); ha il minimo di luce (molto stretta e calore e luce sono limitati), rimuovere la legna è difficile (occorre sempre attraversare un bosco, facendo lievitare i costi);
    \item paralleli alla pendenza: aumenta la luce, ma aumenta il movimento della neve;
    \item diagonalmente alla pendenza: mixa i due vantaggi, l'esposizione dev'essere verso la luce del mattino.
\end{itemize}
I margini regolari della fessura modifica i collettivi, rendendoli maggiormente instabili.\\
Nel piano montano si può fare quello che si vuole: l'abete rosso si presta bene a molti trattamenti:
\begin{itemize}
    \item taglio raso: utilizzato soprattutto il selvicoltura industriale, a nord delle alpi, con tagliate fino 10 ettari, le pendenze sono minime o assenti. In austria le tagliate raso sono limitate a 2-3 ha, con rimboschimenti immediatamente dopo al taglio;
    \item taglio a fratte: in Italia, erano tagliate a raso di circa 3 ha;
    \item taglio a buche, molto utilizzato in Italia e distinguibile in grandi buche (la proiezione a terra delle piante del margine non coprono tutta la buca) e piccole buche (se la buca è più piccola della proiezione delle piante di margine). Con gli attacchi del bostrico, la natura decide la grandezza delle buche, facendo così rinnovare l'abete rosso. Nelle alpi solitamente si lavora con piccole buche (max 500 mq);
    \item  tagli successivi: molto facili, 50\% per la sementazione..... non vengono molto fatti, poichè bisogna tornare molte volte nel cantiere;
    \item tagli a strisce: normalmente la larghezza di 1.5-2 volte l'altezza delle piante del margine e lunghe possibilmente fino a contrafforti robusti;
    \item tagli marginali, simili ai tagli a strisce: si parte da un margine del bosco e si fa una striscia verso un interno del bosco; strutta l'effetto margine (mantello). Il mantello, ecologicamente, è un ecotono: contiene specie tipiche dell'ecotono e degli altri due popolamenti. Se viene tagliato il margine significa che si ha già rinnovazione di altre aree, che vengono così sfruttate;
    \item il difetto dei tagli a strisce è che crea margini stretto e troppo .... andrebbero fatti controvento, in modo che il margine sia protetto;
    \item taglio Femelschlag (tagli successivi perfezionati): taglio giusto, al momento giusto, nel posto giusto. Il risultato è un popolamento coetaneiforme. Si parte da un nucleo di rinnovazione già presente (che viene liberato), facendo un taglio di rinnovazione attorno; la volta dopo che si entra , si svolge una taglio di sgombero sull'area liberata.... In certi posti si fa il taglio di sementazione ed in certi il taglio di sgombero. Molto usato in Alto Adige. Esistono diverse scuole. 
\end{itemize}
Un tipo di pecceta particolare è la pecceta megaforbie: è una pecceta di zone umide, a zone non troppo elevate, ricche di megaforbie. La rinnovazione avviene spesso su legno morto (ceppaia o pianta abbattuta a terra). Spesso ci sono abeti sui trampoli, piante molto stabili. 